% ch3 D.2 自旋波：Heisenberg 哈密顿量与磁状态 (原书页 157–163 / PDF p172–p178)

% 衔接说明：文件最顶部（%--D1-TAIL-- 标记之下）为前一小节 D.1 的收尾文字
%（原书 p.157 顶部、小节 D.2 标题之前的一段），装配时应移至承接原书 p.156
% 的 D.1 文件末尾拼接。

%--D1-TAIL--
因此，把 Wannier 激子当作玻色子来处理是完全合理的。关于由费米子对组成的激发是否可以、以及何时可以当作玻色子来处理这一普遍问题，最详尽的论述见于 Blatt 与 Matsubara 最近的工作；它读起来相当艰涩，但看来确实厘定了这样的事实：这类激发几乎总是可以这样处理。

\subsection{自旋波：Heisenberg 哈密顿量与“磁状态”}

绝缘体中的自旋波不过是 Frenkel 激子的一种修改与推广。它们是这样一种激子：其中未被占据的或“上面的”带并不是 Hartree-Fock 势中一个独立的轨道带，而正是电子由之被激发出来的那同一个带，只不过自旋相反。要理解这一点何以可能，我们必须简要地讨论一下绝缘体中“磁状态”的结构 (91)。

让我们从最简单的情形，即铁磁绝缘体开始；也就是说，我们假定一个 Hartree-Fock 解——不必是最低的那个解，但须自洽——把第 $n$ 个带中所有自旋向上的态全部填满，而所有自旋向下的态全部空着。

\begin{equation*}
\Psi_{\mathrm{HF}} \;=\; \prod_k \bigl(c^{\,n}_{k\uparrow}\bigr)^{\dagger}\, \Psi_{\mathrm{vac}} \;=\; \prod_j \bigl(w^{\,n}_{j\uparrow}\bigr)^{\dagger}\, \Psi_{\mathrm{vac}}
\end{equation*}

其中这两个表达式是等价的（$c$ 是 Bloch 函数的产生算符，$w$ 是 Wannier 函数的产生算符），因为该带是满带。

读者当还记得，在 Frenkel 激子的情形中，我们作为零级近似出发时取了一个原子激发能 $E_0$，它比带激发能 $E_e - E_h$ 小一个量 $U$，这个 $U$ 代表上、下两个带的 Wannier 函数中处于同一原子上的两个电子之间的排斥能。

撇开原子间的相互作用不谈，在磁性情形中 $E_0$ 必须为零，因为被激发的电子进入的正是基态电子原先所在的同一轨道。也就是说，在这一阶段，把一个自旋就地翻转（\emph{in situ}）并不耗费能量，只有把自旋翻转、同时又把电子移到邻近的原子上去才需要能量。若 Wannier 函数为 $a(r-R_j)$，则能量 $U$ 为

\begin{equation*}
U \;=\; \int \frac{e^{2}}{|r_1 - r_2|}\, dr_1\, dr_2\;\, |a(r_1)|^{2}\, |a(r_2)|^{2}
\end{equation*}

显然，磁状态要有任何意义，第一个先决条件就是 $U$ 必须大，特别是相对于“跳跃”即动能积分

\begin{equation*}
b_{jk} \;=\; \int a^{*}(r-R_j)\, \mathcal{H}\, a(r-R_k)\, dr
\end{equation*}

而言要大。

$U$ 必须很大，这一点可以再次通过考虑铁磁情形看出，此时在能量中只计入以下两项：

\begin{equation*}
\mathcal{H}_0 \;=\; \sum_{jk\sigma} b_{jk}\, w^{\dagger}_{j\sigma}\, w_{k\sigma} \;+\; U \sum_j n_{j\uparrow}\, n_{j\downarrow}
\end{equation*}

我们可以在 Hartree-Fock 近似下处理 $H_0$，假定把所有自旋向上的态都填满，使得 $\langle n_{j\uparrow}\rangle = 1$、$\langle n_{j\downarrow}\rangle = 0$。在这个势中，空穴态的能量当然就是

\begin{equation*}
E(k) \;=\; \sum_j\, e^{ik\cdot R_j}\, b_{0j}
\end{equation*}

而电子态的能量则是 $U + E(k)$（图 53）。
%FIG{53}: 图 53　（原书图注仅作 “Figure 53”）电子带与空穴带示意：上方曲线为电子带 $U+E(k)$，下方曲线为空穴带 $E(k)$，两曲线极小值之间的竖直箭头标出间隔 $U$
这种局面显然只有在所有电子态都位于所有空穴态之上时才是稳定的，即

\begin{equation*}
\bigl[E(k)\bigr]_{\max} \;<\; U \;+\; \bigl[E(k)\bigr]_{\min}
\end{equation*}

这是电子不至于自发地从自旋向上的态转入自旋向下的态、从而把铁磁绝缘体变成非磁性金属的必要条件，但远非充分条件。

现在我们可以从第二个角度来研究铁磁绝缘体的稳定性。这就是说，我们可以计算“激子”或“自旋激发波”的能量——把一个电子从自旋向上态激发到自旋向下态、并让由此形成的束缚复合体在晶格中传播。我们将发现，$k=0$ 的激子——让我们使用“自旋波”这个术语——总是具有零能量；但有限波数的自旋波却可能具有正的或负的能量。如果这些能量是负的，这也代表铁磁系统的一种失稳，大概是朝自旋重新取向、改排成反铁磁阵列的方向。

处理这一问题的这条途径是 Slater (92) 以及 Schubin 与 Vonsovsky (93) 的途径。它虽然在历史上是第一个真正的多体处理，却并未阐明从这一系统所能获得的全部信息。更为完全、而在严格性上本质上等价的，是“矢量模型”（vector model）的处理方式，我们稍后即来讨论。

Coulomb 相互作用不含有任何直接的偶极矩阵元，因为自旋翻转跃迁受到自旋与宇称的双重禁戒。这样一来，激发只能借助所谓的“交换”矩阵元从一个原子传递到另一个原子：

\begin{equation*}
J_{jk} \;=\; \int a^{*}(r-R_j)\, a(r-R_k)\, \frac{e^{2}}{|r-r'|}\, a^{*}(r'-R_k)\, a(r'-R_j)\, dr\, dr'
\end{equation*}

它就是“交叠电荷密度”（overlap charge density）

\begin{equation*}
\rho_{jk} \;=\; a^{*}(r-R_j)\, a(r-R_k)
\end{equation*}

的 Coulomb 自能。注意，因为 $J$ 是一个 Coulomb 自能，它必定为正。

$J$ 与费米子算符的两种不同编组相乘：

\begin{align*}
\mathcal{H}_{\text{direct exchange}} \;=\; \frac{1}{2} \sum_{jk} J_{jk}\, \Bigl\{ &\Bigl[\, w^{\dagger}_{j\uparrow} w^{\dagger}_{k\uparrow} w_{j\uparrow} w_{k\uparrow} \\
&+\, w^{\dagger}_{j\downarrow} w^{\dagger}_{k\downarrow} w_{j\downarrow} w_{k\downarrow} \Bigr] \;+\; \Bigl[\, w^{\dagger}_{j\uparrow} w^{\dagger}_{k\downarrow} w_{j\uparrow} w_{k\downarrow} \\
&+\, w^{\dagger}_{j\downarrow} w^{\dagger}_{k\uparrow} w_{j\downarrow} w_{k\uparrow} \Bigr] \Bigr\}
\end{align*}

前一对项可以改写为

\begin{align*}
&-\bigl(n_{j\uparrow}\, n_{k\uparrow} \;+\; n_{j\downarrow}\, n_{k\downarrow}\bigr) \\
&=\; -\frac{1}{2}\,\Bigl\{ \bigl(n_{j\uparrow} + n_{j\downarrow}\bigr)\bigl(n_{k\uparrow} + n_{k\downarrow}\bigr) \;+\; \bigl(n_{j\uparrow} - n_{j\downarrow}\bigr)\bigl(n_{k\uparrow} - n_{k\downarrow}\bigr) \Bigr\} \\
&=\; -\frac{1}{2}\,\bigl(1 + \sigma_{zj}\, \sigma_{zk}\bigr)
\end{align*}

这里假定了每个原子上总有一个电子：

\begin{equation*}
n_{j\uparrow} + n_{j\downarrow} \;=\; 1
\end{equation*}

我们已经把这一对费米子认同于一个旋量（spinor）；在这一情形中，这正是我们通常所说的与给定原子相联系的自旋算符。

后一对项则是更为常见的激子转移项；我们看到

\begin{equation*}
w^{\dagger}_{j\uparrow}\, w_{j\downarrow} \;=\; \sigma_{j+} \;=\; \frac{1}{2}\left(\sigma_{jx} + \sigma_{jy}\right)
\end{equation*}

等等，于是总的交换能量可以用自旋算符写成

\begin{align*}
\mathcal{H}_{\text{direct exchange}} \;&=\; -\frac{1}{4} \sum_{jk} J_{jk}\left(1 + \sigma_{zj}\sigma_{zk} + \sigma_{xj}\sigma_{xk} + \sigma_{yj}\sigma_{yk}\right) \\
&=\; -\frac{1}{4} \sum_{jk} J_{jk}\left(1 + \boldsymbol{\sigma}_{j}\cdot\boldsymbol{\sigma}_{k}\right) \tag{23}
\end{align*}

自旋记号的优点在于，它至少在形式上使我们不必把思路局限于研究单个自旋波，而可以考虑样品中任意多个自旋同时反转的情形。然而它依赖于一个基本假定：$U$ 非常大，到所有其他带的激发能也非常大，因而我们可以取 $n_{j\uparrow} + n_{j\downarrow} = 1$。当然，这只有在 $U\to\infty$ 这个非物理的极限下才严格成立；但正如 Frenkel 激子的情形一样，我们可以求助于投影到使其严格成立的态流形之上的微扰论。只要 $b/U$ 足够小，这样的微扰论看来无疑会收敛［尽管 Slater 曾对这一点提出质疑，并没有给出令人信服的理由 (94)］。

另一方面，只要我们把激发彼此之间相互作用的结果严格用自旋 $\boldsymbol{\sigma}$ 表出，就没有任何东西阻止我们处理存在不定数目的自旋反转的情形。这样一来，这就实在是一种相当新型的元激发理论：我们允许不定程度的激发，$n_{\text{exc}}\sim N$，但只允许某种特殊类型的激发。对于新理论的有效哈密顿量 (23)，我们也许能够精确处理，也许不能；但只要我们拒绝允许真实——有别于虚的——激发中的电子永久地离开其相应空穴所在的格点，它就是一个准精确的哈密顿量，其意义几乎与有效质量理论相同。

在这一理论中，自旋波的运动与相互作用的第二种、也是十分重要的机制，就是我们在激子情形中讨论过的同一个微扰论项。一个自旋反转的电子可以虚跳跃到邻近的原子上去，从而赋予自旋波一个有限大小的波函数。我们可以直接照搬在激子情形中得到的结果。二级能量为（这里计及 $b_{j\uparrow k\uparrow} = b_{j\downarrow k\downarrow}$）

\begin{align*}
\mathcal{H}^{(2)} \;&=\; \frac{1}{U} \sum_{jk} |b_{jk}|^{2}\, \Bigl\{ n_{j\uparrow}\, n_{k\uparrow} \;+\; n_{j\downarrow}\, n_{k\downarrow} \\
&\qquad\quad +\; w_{j\uparrow} w_{k\downarrow} w_{k\uparrow} w_{j\downarrow} \;+\; w_{j\downarrow} w_{k\uparrow} w_{k\downarrow} w_{j\uparrow} \Bigr\} \\
&=\; \frac{1}{2} \sum_{jk} \frac{|b_{jk}|^{2}}{U}\left(\sigma_{zj}\sigma_{zk} + \sigma_{xj}\sigma_{xk} + \sigma_{yj}\sigma_{yk} - 1\right) \\
&=\; \frac{1}{2} \sum_{jk} \frac{|b_{jk}|^{2}}{U}\left(\boldsymbol{\sigma}_{j}\cdot\boldsymbol{\sigma}_{k} - 1\right)
\end{align*}

无需对自旋波问题作任何认真的研究就可以明显看出：若 $\frac{1}{2}J_{jk} > |b_{jk}|^{2}/U$，铁磁态将是稳定的，因为此时让所有的 $\boldsymbol{\sigma}_j\cdot\boldsymbol{\sigma}_k$ 尽可能为正便可使能量取极小；而在相反的情形下，某种反铁磁态大概会更稳定。请允许我离题片刻，谈谈绝缘体中交换相互作用里这两个最重要项的物理来源。$J_{jk}$ 这个项相当直接地就是两个正交轨道 $a(r-R_j)$ 与 $a(r-R_k)$ 中自旋平行的电子在相互作用能中的 Hartree-Fock 交换项。因为两个电子自旋平行，在它们波函数交叠的区域内存在一个“交换空穴”（exchange hole），使它们平行时能量降低。这一能量收益正是原子中 Hund 定则的来源，该定则说：同一壳层中电子的自旋在可能时将相互平行。

这个交换项必然为正——那么，围绕交换积分的符号而激烈进行的那些争论又是怎么回事呢？原因在于：两个原子相耦合的 Heitler-London 理论给出的交换积分并不只是这一项，而是一个涉及完整哈密顿量、并用不必正交的任意函数代替 $a(r-R_j)$ 的积分。这样的积分可以取任意符号，而且实际上通常为负；最常被计算的正是这个 H-L 积分。虽然 Herring (95) 已经证明，对相距相当远的单电荷原子，Heitler-London 程序几乎严格正确，但在我看来，把主要的反铁磁效应作为一个单独的项 $\sim -b^{2}/U$ 计入，会使物理更为清晰。之所以出现这一项，是因为当两个电子的自旋反平行时，不相容原理不再要求它们具有严格正交的波函数；它们可以因溢入彼此的波函数而获得动能。

可以证明，当两个波函数并非因对称性而自动正交时，$b^{2}/U$ 项几乎总是占主导地位；例如，根据 Wigner (95) 的一个定理，任意普通势场中两个电子的\emph{最低}态总是单态（singlet），表明交换是反铁磁性的。

但是不相容原理的作用常常——比如在原子中——使得两个电子的波函数自动正交，从而 $b_{jk}$ 为零。服从这一原理的绝缘铁磁体的若干孤立例子，如 $\mathrm{CrO_2}$，如今已经知道；但数以百计的其他例子表现出的是反铁磁交换。

交换效应重要的那些常见磁性材料，是铁族金属——偶尔也有其他过渡族金属——的氧化物或其他相当浓的盐类。一个典型的例子是 $\mathrm{NiO}$。在这种材料中，我们感兴趣的带主要由 Ni 离子上的 $d$ 态构成，不过由于 O 的 $2p$ 电子与 Ni 的 $d$ 态之间的共价键合，存在相当大的混杂。由文献 (91) 可知，$b$ 的合理数值约为 1 ev，$U$ 约为 10 ev，$J$ 约为 $b^{2}/U$ 的 1/3。

在这种通常的情形，即 $2(|b_{jk}|^{2}/U) > J_{jk}$ 的情形中，相互作用被称为“反铁磁的”，电子自旋在基态或在 $T=0$ 附近将倾向于使 $\boldsymbol{\sigma}_j\cdot\boldsymbol{\sigma}_k$ 尽可能为负。观测到的态往往相当好地近似于由两个或多个交替自旋子晶格构成的阵列（图 54）。当然，在这样的情形中，我们所处理的确实是这样一种局面：
%FIG{54}: 图 54　（原书图注仅作 “Figure 54”）交替自旋的子晶格阵列示意：三行格点各带向上、向下交替排列的自旋箭头；最下一行在格点 $j$、$k$ 处以弧形箭头画出一对正在交换的自旋
从铁磁的观点看，它含有非常高度的激发。然而正如我们所指出的，只要我们相信微扰论从 $n_\uparrow + n_\downarrow = 1$ 的准简并态流形出发是有效的、并且不允许向自由电子态作真实（有别于虚的）跃迁，这是完全没有问题的。

这样一种理论的有效性判据必须是：$b$ 应当比 $U$ 小得多。读者当还记得，在 Frenkel 激子情形中，当 $b$ 相对 $U$ 变大时，我们有一个从 Frenkel 型到 Wannier 型激子的快速的、但大概并非不连续的转变；而在这里，我们可以容易地看出，这个转变将是灾难性的。

正如同 Frenkel 激子的情形一样，关键问题围绕激子的结合能。然而在这里，我们并没有那个把电子从同一原子上的 $s$ 态激发到 $p$ 态所需的附加激发能 $E_0$；唯一阻止系统退化为纯粹金属态的能量，是电子与各自空穴之间的结合能。把电子从它的空穴中解放出来，我们可以获得量级为 $\sim Zb$ 的能量，但我们失去能量 $U$，而允许电子虚离开其空穴（以及相反的过程）所获得的只是二级能量 $-Zb^{2}/U$。于是我们可以估计 $Zb \sim U$ 作为从反铁磁体转变到金属的转变点的合理数值。此外，当 $b$ 大时，介电屏蔽再度成为问题，$U$ 自然也就更小。

处理向金属态转变这一问题的另一种、同样相当近似的方式，是问交替阵列究竟是不是多体问题的一个稳定的 Hartree 解。也就是说，我们可以问：$k$ 原子处自旋向下的电子所感受到的 Hartree 场相对于 $j$ 原子处的如何。名义上，这个差值是 $U$，即 $j$ 处自旋向上的电子与自旋向下的电子之间的排斥能。但在 Hartree 理论中，我们应当计及这样一个事实：$j$ 处的自旋向上电子并非时时刻刻都真正位于 $j$ 处，所以 Hartree 场的实际差值可能要小得多。

格点 $j$ 上的自旋向上电子以概率 $|b_{jk}/U|^{2}$ 跳到它的每一个近邻格点 $k$；同样，格点 $k$ 也可能以概率 $|b_{jk}/U|^{2}$ 被一个自旋向上电子占据。于是，$j$ 处的自旋向下态相对于 $k$ 处的自旋向下态的实际能量差为

\begin{align*}
\Delta E \;&=\; U\left(\langle n_{j\uparrow}\rangle - \langle n_{k\uparrow}\rangle\right) \\
&=\; U\left(1 - 2Z\,\Bigl|\frac{b_{jk}}{U}\Bigr|^{2}\right)
\end{align*}

但实际上，决定混杂大小的能量分母不应当是 $U$ 本身，而应当是 $\Delta E$，于是我们得到一个自洽方程

\begin{equation*}
\Delta E \;=\; U\left(1 - \frac{2}{(\Delta E)^{2}} \sum_{k} |b_{jk}|^{2}\right)
\end{equation*}

或

\begin{equation*}
(\Delta E)^{3} \;-\; U\,(\Delta E)^{2} \;=\; -\,2UZ\,|b_{jk}|^{2}
\end{equation*}

这个方程可以像图 55 所示那样图示出来。于是 $b$ 的最大可能值由下式给出：

\begin{equation*}
\left(\frac{b}{U}\right)_{\max} \;\cong\; \left(\frac{2}{27Z}\right)^{1/2} \;\sim\; \frac{1}{10}\ \text{左右}
\end{equation*}

在任一高于这个临界值的 $b$ 下，磁状态再也无法维持。它将对电子的重新分布不稳定：重新分布使两种自旋态的占据相等，态成为金属型的。铁系中较低的氧化物，

% ——上句在原书 p.163 末中断于 “The lower iron-series oxides,”，p.164（PDF p179）起由后续代理续译衔接。
